\documentclass{article}

\PassOptionsToPackage{numbers, sort, compress}{natbib}
\usepackage[main,final]{neurips_2025}

\usepackage[T1]{fontenc}
\usepackage[utf8]{inputenc}
\usepackage{hyperref}
\usepackage{url}
\usepackage{booktabs}
\usepackage{nicefrac}
\usepackage{microtype}
\usepackage{xcolor}

\usepackage{subcaption}
\usepackage{graphicx}
\usepackage{multirow}
\usepackage{amsmath,amssymb,amsfonts}
\usepackage{amsthm}
\usepackage{mathrsfs}
\usepackage{textcomp}
\usepackage{manyfoot}
\usepackage{algorithm}
\usepackage{algorithmicx}
\usepackage{algpseudocode}
\usepackage{listings}

\newtheorem{theorem}{Theorem}
\newtheorem{proposition}[theorem]{Proposition}
\newtheorem{lemma}{Lemma}
\newtheorem{example}{Example}
\newtheorem{remark}{Remark}
\newtheorem{definition}{Definition}
\newtheorem{assumption}{Assumption}


\title{PersonGen: Knowledge-Preserving Personalization of Text-to-Image Diffusion Models}


\author{
  Karina Peskova\\
  Lomonosov Moscow State University\\
  Moscow, Russia\\
  \texttt{email@example.com}\\
}


\begin{document}

\maketitle


\begin{abstract}

Recent text-to-image diffusion models demonstrate remarkable generative
capabilities but require adaptation to generate personalized concepts.
Fine-tuning approaches such as DreamBooth and LoRA allow efficient
personalization, but may alter previously learned knowledge and reduce
generalization ability.

In this work, we study the problem of knowledge-preserving personalization
of text-to-image diffusion models. We combine LoRA-based personalization
with AlphaEdit-inspired model editing methods and investigate how personalized
concept information is represented inside diffusion models.

We analyze the contribution of cross-attention key/value parameters to concept
transfer and study whether editing these components is sufficient for preserving
the original model capabilities while introducing new visual concepts.

\end{abstract}


\textbf{Keywords:} text-to-image diffusion models, personalization, LoRA,
DreamBooth, model editing, Stable Diffusion, knowledge preservation


\section{Introduction}\label{sec:intro}

Text-to-image diffusion models have recently achieved impressive results in
high-quality image generation. However, adapting these models to new concepts,
objects, or subjects remains a challenging problem. A personalized model should
learn the target concept while preserving the knowledge acquired during
pretraining.

Existing personalization methods, such as DreamBooth and LoRA, efficiently
adapt pretrained diffusion models using a small number of images. However,
fine-tuning can introduce unintended modifications and reduce the ability of
the model to generate previously learned concepts.

An alternative direction is model editing, where model parameters are modified
in a constrained manner to introduce new knowledge while minimizing changes to
the original model. This work investigates whether such approaches can be
combined with personalization methods for more controlled concept transfer.

The main goal of this research is to analyze how personalized concepts are
stored in text-to-image diffusion models and whether knowledge-preserving
editing methods can improve personalization quality.

\textbf{Contributions.} Our contributions can be summarized as follows:

\begin{itemize}
    \item We develop a pipeline combining DreamBooth LoRA personalization with
    AlphaEdit-based model editing for Stable Diffusion models.
    
    \item We analyze the role of cross-attention key/value matrices in storing
    personalized concepts.
    
    \item We perform ablation studies over textual identifiers, LoRA training
    parameters, checkpoints, and editing configurations.
    
    \item We investigate the trade-off between concept strength and preservation
    of pretrained model knowledge.
\end{itemize}


\textbf{Outline.}
The rest of the paper is organized as follows.
Section~\ref{sec:rw} describes related work.
Section~\ref{sec:prelim} introduces preliminary concepts.
Section~\ref{sec:method} presents the proposed approach.
Section~\ref{sec:exp} describes experimental evaluation.
Section~\ref{sec:disc} discusses the obtained results.


\section{Related Work}\label{sec:rw}


\textbf{Text-to-image personalization.}

Personalization methods adapt pretrained diffusion models to generate specific
subjects or concepts using a limited number of examples. DreamBooth introduces
subject-driven generation by fine-tuning model parameters, while LoRA provides
a parameter-efficient alternative by learning low-rank updates.


\textbf{Knowledge-preserving model editing.}

Model editing methods aim to modify neural networks while minimizing unwanted
changes to existing knowledge. AlphaEdit introduces constrained editing
strategies that project parameter updates into subspaces preserving previous
model behavior.


\textbf{Concept representation in diffusion models.}

Previous studies analyze how textual representations and attention mechanisms
affect image generation. In this work, we focus on cross-attention key/value
matrices and their role in transferring personalized concepts.


\section{Preliminaries}\label{sec:prelim}


\subsection{General notation}

Let $M_0$ denote a pretrained Stable Diffusion model and $M_p$ denote a
personalized model obtained after fine-tuning.

The LoRA personalization process can be represented as:

\[
W_{LoRA}=W_0+\Delta W ,
\]

where $W_0$ is the original parameter matrix and $\Delta W$ is a low-rank
adaptation.


\subsection{Assumptions}

We assume that:

\begin{itemize}
    \item the pretrained diffusion model contains general visual knowledge;
    \item personalization modifies only a subset of model parameters;
    \item cross-attention layers play an important role in connecting textual
    representations with visual concepts.
\end{itemize}


\section{Method}\label{sec:method}


The proposed pipeline consists of several stages.

\subsection{LoRA personalization}

First, DreamBooth LoRA fine-tuning is performed to learn a new visual concept
from a small image dataset. A special identifier token is introduced to represent
the personalized concept.


\subsection{Attention parameter extraction}

After personalization, cross-attention key and value projection matrices are
extracted from the diffusion model.

These matrices are analyzed as possible carriers of personalized information.


\subsection{Text representation analysis}

We construct text embedding matrices corresponding to:

\begin{itemize}
    \item identifier tokens;
    \item content words;
    \item all valid tokens.
\end{itemize}


\subsection{AlphaEdit-based modification}

The extracted representations are used to compute constrained parameter updates.
The goal is to transfer the personalized concept while minimizing changes to
the original model.


\section{Experiments}\label{sec:exp}


We conduct experiments using Stable Diffusion v1-4.

The experimental evaluation includes:

\begin{itemize}
    \item comparison of different identifier tokens;
    \item LoRA checkpoint analysis;
    \item learning rate ablations;
    \item comparison of different text representation strategies;
    \item interpolation between LoRA and edited weights.
\end{itemize}


For interpolation experiments we define:

\[
W(\alpha)=(1-\alpha)W_{edit}+\alpha W_{LoRA},
\]

where $\alpha=0$ corresponds to the edited model and $\alpha=1$ corresponds
to the original LoRA personalization.


Evaluation is performed using generated image comparisons and analysis of
concept preservation.


\section{Discussion}\label{sec:disc}


The experiments demonstrate that personalized information is distributed across
multiple components of diffusion models.

Although modifying cross-attention parameters can transfer part of the concept,
the full personalization behavior cannot always be recovered using only
key/value matrices.

This suggests that successful personalization requires considering a broader set
of parameters responsible for concept representation.


\section{Conclusion}\label{sec:concl}


In this work, we investigate knowledge-preserving personalization of
text-to-image diffusion models.

We propose a research pipeline combining DreamBooth LoRA personalization and
AlphaEdit-based editing. Our analysis provides insights into the localization
of personalized knowledge inside diffusion models and the challenges of
performing controlled concept transfer.


%%%%%%%%%%%%%%%%%%%%%%%%%%%%%%%%%%%%%%%%%%%%%%%%%%%%%%%%%%%%

\bibliographystyle{unsrtnat}
\bibliography{references}


\newpage
\appendix

\section{Appendix / supplemental material}\label{app}


\subsection{Additional experiments / Proofs of Theorems}\label{app:exp}


Additional experimental results, generated image grids, and ablation studies
are provided in the supplementary material.


\end{document}